\section{XÂY DỰNG TẬP DỮ LIỆU VÀ MÔ PHỎNG SUY GIẢM}
\subsection{Giới thiệu bộ dữ liệu GTSRB}
German Traffic Sign Recognition Benchmark (GTSRB) là bộ dữ liệu phân loại biển báo giao thông với 43 lớp~\cite{gtsrb}. Trong dữ liệu dự án, tập huấn luyện chính thức có 39.209 ảnh và tập kiểm thử chính thức có 12.630 ảnh. Các ảnh khác nhau về góc quan sát, kích thước biển báo, nền cảnh và điều kiện chiếu sáng. Đây là những yếu tố khiến việc tăng cường ảnh cần giữ được đặc trưng hình học và màu sắc thay vì chỉ nâng giá trị sáng trung bình.

Một số ảnh huấn luyện được chụp liên tiếp từ cùng một biển báo. Những ảnh thuộc cùng chuỗi quan sát có mức tương đồng cao. Nếu chia train và validation ngẫu nhiên ở cấp ảnh, các khung gần giống nhau có thể xuất hiện ở cả hai tập, làm kết quả validation lạc quan hơn khả năng tổng quát thực tế. Vì vậy, dự án chia dữ liệu theo lớp và \term{track}, giữ các ảnh cùng track trong một tập.

\begin{table}[htbp]
\centering
\caption{Phân chia dữ liệu trong quy trình ResNet-18.}
\begin{tabular}{P{4.0cm}rP{7.3cm}}
\toprule
Tập dữ liệu & Số ảnh & Vai trò \\
\midrule
Huấn luyện & 31.350 & Huấn luyện ResNet-18 trên ảnh gốc \\
Validation & 7.829 & Chọn checkpoint và đánh giá trong huấn luyện \\
Cách ly & 30 & Loại track có ảnh trùng pixel với test khỏi huấn luyện \\
Official test & 12.630 & Kiểm thử toàn bộ tập chính thức \\
Validation chọn xử lý & 1.290 & 30 ảnh/lớp để chọn tham số tăng cường \\
Test khảo sát & 430 & 10 ảnh/lớp cho notebook và ablation \\
\bottomrule
\end{tabular}
\label{tab:splits}
\end{table}

Tập cách ly gồm 30 ảnh thuộc track có ảnh trùng pixel với tập test, theo kiểm kê được ghi trong dự án. Toàn bộ track được tách ra để tránh chỉ loại những ảnh trùng chính xác nhưng vẫn giữ các khung gần kề. Tổng số ảnh train, validation và cách ly bằng 39.209; test chính thức được giữ riêng. Hạt giống chia dữ liệu là 42.

Các tập 1.290 và 430 ảnh là những tập con của validation và test, không phải dữ liệu bổ sung. Việc lấy cùng số ảnh mỗi lớp giúp khảo sát cân bằng 43 lớp và giảm chi phí chạy nhiều phương pháp. Tuy nhiên, kết quả trên 430 ảnh không đại diện đầy đủ cho phân bố tần suất của 12.630 ảnh test; báo cáo vì thế trình bày hai phạm vi này thành các bảng riêng.

\subsection{Tiền xử lý và chuẩn hóa ảnh}
\subsubsection{Nguồn nhãn và hình thức học}
Gói \path{GTSRB_Final_Training_Images.zip} đã chứa ảnh và annotations; mỗi thư mục lớp có file \texttt{GT-<ClassID>.csv}. Cột \texttt{ClassId} cung cấp nhãn 0--42, các cột ROI cung cấp vị trí biển báo. Nhãn test được tải riêng trong \path{GTSRB_Final_Test_GT.zip}~\cite{gtsrbdata}. Vì vậy, không tải một file GT train riêng không đồng nghĩa với dữ liệu train không có nhãn.

Bộ phân loại ResNet-18 học có giám sát: mỗi mẫu trả về cặp ảnh--nhãn, dùng cross-entropy và cập nhật toàn mạng. Ba enhancer học sâu chỉ nạp trọng số tác giả để suy luận; Zero-DCE gốc có cơ chế zero-reference, nhưng điều đó không làm bài toán phân loại GTSRB của dự án trở thành học không giám sát. Các thuật toán truyền thống không có bước huấn luyện.

\subsubsection{Cắt vùng quan tâm}
Vùng quan tâm (Region of Interest, ROI) được xác định bằng tọa độ chú thích của dữ liệu. Dự án mở rộng vùng này thêm 5\% kích thước theo mỗi phía, đồng thời giới hạn khung cắt trong kích thước ảnh. Phần lề nhỏ giúp giữ đường viền biển báo và một lượng ngữ cảnh xung quanh; nếu cắt sát quá mức, những pixel biên có thể bị loại bỏ trước khi thuật toán tăng cường hoạt động.

Với ROI $(x_1,y_1,x_2,y_2)$ có điểm cuối bao gồm trong vùng, mã nguồn chuyển $x_2,y_2$ sang mốc cuối không bao gồm bằng cách cộng 1. Sau đó tính độ mở rộng theo chiều rộng và chiều cao, làm tròn ra ngoài bằng phép lấy sàn/lấy trần và chặn tại biên ảnh. Quy ước này tránh bỏ mất hàng hoặc cột pixel cuối cùng khi cắt bằng thư viện ảnh.

\subsubsection{Đổi kích thước và miền cường độ}
Ảnh ROI được đổi sang RGB, nội suy song tuyến tính về $64\times64$ rồi lưu dưới dạng chuẩn hóa:
\begin{equation}
X\in[0,1]^{64\times64\times3},\qquad X=\frac{I_{8\text{-bit}}}{255}.
\end{equation}
Nội suy song tuyến tính sử dụng tổ hợp có trọng số của các pixel lân cận. Cách làm này phù hợp với yêu cầu đưa các vùng biển báo có kích thước khác nhau về cùng đầu vào, nhưng có thể làm mềm các nét vốn rất nhỏ. Do đó, quá trình phóng lên $224\times224$ sau tăng cường không tạo lại thông tin đã mất ở bước chuẩn hóa $64\times64$.

Các thuật toán làm việc với số thực \texttt{float32}. Riêng HE và CLAHE chuyển ảnh sang số nguyên 8 bit để dùng các phép xử lý histogram của OpenCV rồi chuyển trở lại số thực. Việc lượng tử hóa ở hai phương pháp này là một phần của cách cài đặt, cần được phân biệt với các phương pháp thực hiện trực tiếp trên giá trị liên tục.

\subsection{Mô hình suy giảm thiếu sáng}
Mô hình suy giảm được áp dụng theo thứ tự cố định: Gamma làm tối, nhiễu Poisson, nhiễu Gaussian và cuối cùng chặn về $[0,1]$. Với một pixel của ảnh gốc $X$, quá trình được mô tả bởi:
\begin{align}
U &= X^{\gamma_d},\label{eq:dark}\\
V &= \frac{\mathcal{P}(pU)}{p},\label{eq:poisson}\\
Y &= \operatorname{clip}\bigl(V+\varepsilon,0,1\bigr),\quad
\varepsilon\sim\mathcal{N}(0,\sigma_n^2).\label{eq:degradation}
\end{align}
Trong đó $\gamma_d>1$ điều khiển mức tối, $p$ là hệ số Poisson và $\sigma_n$ là độ lệch chuẩn nhiễu Gaussian. Phép lũy thừa áp dụng theo từng kênh RGB trong miền ảnh chuẩn hóa; đây là mô hình suy giảm tổng hợp, không phải mô phỏng đầy đủ quá trình thu nhận ánh sáng của cảm biến RAW.

\subsubsection{Vai trò của Gamma làm tối}
Với $0<X<1$ và $\gamma_d>1$, ta có $X^{\gamma_d}<X$. Khi tăng $\gamma_d$, vùng trung gian bị kéo về cường độ thấp mạnh hơn. Ví dụ, một pixel có cường độ 0,5 trở thành khoảng 0,354; 0,218 và 0,125 khi $\gamma_d$ lần lượt bằng 1,5; 2,2 và 3,0. Điểm đen 0 và điểm trắng 1 vẫn được giữ ở hai đầu miền.

Điều này giải thích vì sao ảnh tối có histogram dịch về bên trái mà vùng phản xạ rất sáng vẫn có thể còn rõ. Phép Gamma tạo được nhiều mức suy giảm với tham số dễ kiểm soát, đồng thời vẫn giữ quan hệ xác định giữa ảnh gốc và ảnh tối khi chưa thêm nhiễu.

\subsubsection{Vai trò của nhiễu Poisson và Gaussian}
Trước bước chặn, nhiễu Poisson trong công thức~\eqref{eq:poisson} có kỳ vọng $U$ và phương sai $U/p$. Khi $p$ giảm, mức dao động tăng. Mặc dù phương sai tuyệt đối nhỏ ở vùng tối, độ lệch chuẩn tương đối so với tín hiệu là $1/\sqrt{pU}$, có thể lớn khi $U$ rất nhỏ. Đây là đặc điểm quan trọng khi khảo sát tính nhạy nhiễu của phép tăng sáng.

Nhiễu Gaussian được cộng với độ lệch chuẩn cố định cho mọi pixel. Trong mô hình này, nó đại diện cho một thành phần nhiễu cộng đơn giản. Khi tín hiệu sau Gamma nhỏ, cùng một mức nhiễu Gaussian sẽ gây ảnh hưởng tương đối lớn hơn đến nội dung. Bước chặn cuối cùng bảo đảm ảnh nằm trong miền hợp lệ, nhưng cũng làm phân bố nhiễu gần 0 hoặc 1 không còn hoàn toàn đối xứng.

\begin{table}[htbp]
\centering
\caption{Ba mức suy giảm sử dụng trong thí nghiệm.}
\begin{tabular}{lrrr}
\toprule
Mức độ & $\gamma_d$ & Poisson $p$ & Gaussian $\sigma_n$ \\
\midrule
Nhẹ (mild) & 1,5 & 120 & 0,005 \\
Vừa (medium) & 2,2 & 60 & 0,010 \\
Nặng (severe) & 3,0 & 30 & 0,020 \\
\bottomrule
\end{tabular}
\label{tab:degrade}
\end{table}

Ba mức trong Bảng~\ref{tab:degrade} thay đổi đồng thời độ tối và hai thành phần nhiễu. Vì vậy, sự chênh lệch giữa mức nhẹ và mức nặng phản ánh tổng tác động của cả ba yếu tố. Muốn xác định riêng ảnh hưởng của nhiễu, cần dùng thí nghiệm ablation giữ nguyên Gamma và lần lượt bật/tắt từng loại nhiễu, được trình bày ở phần thảo luận.

\subsection{Khảo sát histogram và cấu trúc ảnh}
\fig{degradation.png}{1}{Ảnh tham chiếu và ba mức suy giảm trên cùng một crop validation. Hàng dưới là histogram cường độ xám; trục đứng biểu thị số pixel. Hình trích trực tiếp từ notebook ResNet-18, noise seed 101.}{fig:degradation}
Quan sát Hình~\ref{fig:degradation}, vùng nền sáng của biển báo suy giảm rõ khi tăng mức tối. Các nét số trở nên khó phân biệt với nền và nhiễu màu xuất hiện trên cả vùng tương đối đồng nhất. Histogram tập trung ngày càng nhiều tại cường độ thấp; độ tương phản giữa các cấu trúc hữu ích không thể được suy ra chỉ từ việc histogram còn trải rộng đến gần 1.

Trên riêng ảnh minh họa này, PSNR của mức nhẹ, vừa và nặng lần lượt là 17,2577; 12,0942; 9,5306 dB. SSIM tương ứng là 0,7120; 0,5494; 0,3855. Đây là số đo của một ảnh validation được chọn theo quy tắc trong notebook, không phải giá trị trung bình của toàn bộ tập test.

Histogram chỉ biểu diễn tần suất cường độ, không chứa vị trí của pixel. Hai ảnh có histogram gần nhau vẫn có thể khác hoàn toàn về bố trí đường viền hoặc chữ số. Vì vậy, phân tích histogram trong báo cáo luôn đi kèm quan sát ảnh, SSIM và kết quả nhận dạng, tránh kết luận chất lượng chỉ từ sự dịch chuyển phân bố mức xám.

\subsection{Bảo đảm điều kiện so sánh công bằng}
Hạt giống sinh nhiễu được xác định từ noise seed, mã mẫu và tên điều kiện. Ba noise seed là 101, 202 và 303. Với cùng mã ảnh, cùng mức suy giảm và cùng noise seed, tất cả các phương pháp nhận đúng một ảnh tối. Cách thiết kế này tránh việc một thuật toán vô tình được đánh giá trên nhiễu nhẹ hơn thuật toán khác.

Ảnh gốc chỉ được sử dụng để tính chất lượng sau tăng cường. Nhãn lớp chỉ phục vụ huấn luyện bộ nhận dạng, chọn tham số bằng validation và tính metric. Các tập con, cấu hình và checkpoint được khóa trước test. Những ràng buộc này giúp quy kết thay đổi kết quả chủ yếu cho phép xử lý ảnh đang khảo sát trong phạm vi giao thức đã định.
